\documentclass[11pt]{amsart}
\usepackage{geometry}                % See geometry.pdf to learn the layout options. There are lots.
\geometry{letterpaper}                   % ... or a4paper or a5paper or ... 
%\geometry{landscape}                % Activate for for rotated page geometry
%\usepackage[parfill]{parskip}    % Activate to begin paragraphs with an empty line rather than an indent
\usepackage{graphicx}
\usepackage{amssymb}
\usepackage{epstopdf}
\usepackage{multirow}
\DeclareGraphicsRule{.tif}{png}{.png}{`convert #1 `dirname #1`/`basename #1 .tif`.png}

\title{Problema 2.2}

%\date{}                                           % Activate to display a given date or no date

\begin{document}
\maketitle
%\section{}
%\subsection{}
\noindent Construya un diagrama de flujo tal que dado como datos los valores enteros $P$ y $Q$, determine si los mismos satisfacen la siguiente expresión:

\[P^5 + Q^4 - 2 \cdot P^2 < 680\]

\noindent En caso afirmativo, debe imprimir los valores $P$ y $Q$.

\noindent Datos: $P$,$Q$ (variables de tipo entero que expresan los datos que se ingresan).\\

\noindent \textbf {Explicación de las variables} \\

\noindent P,Q: Variables de tipo entero.

\noindent EXP: Variable de tipo real. Almacena el resultado del cálculo de la expresión.\\

\noindent A continuación en la tabla 2.9 podemos observar el seguimiento del algoritmo.

\begin{table}[htbp]
\centering
\label{tab:corrida}
\begin{tabular}{|c|c|c|c|c|}
\hline
\multicolumn{5}{ |c| }{\textbf{Tabla 2.16}} \\ \hline
\multicolumn{2}{|c|}{\multirow{2}{*}{\textbf{NUMERO DE CORRIDA}}} & \textbf{DATOS} & {\multirow{2}{*}{\textbf{CALCULO AUXILIAR}}} & \textbf {RESULTADO}  \\
\multicolumn{2}{|c|}{} &  \textbf{P} \quad  \textbf{Q}  &   &  \textbf{P} \quad \quad \textbf{Q}  \\
\hline
\multicolumn{2}{|c|}{1} &  3  \quad  5 &  634 &  3 \quad \quad 5\\
\hline
\multicolumn{2}{|c|}{2} &  6  \quad  8 &  4240 & \\
\hline
\multicolumn{2}{|c|}{3} &  2  \quad  4 &  256 &  2 \quad \quad 4\\
\hline
\multicolumn{2}{|c|}{4} &  7  \quad  5 &  870 &  \\
\hline
\multicolumn{2}{|c|}{5} &  2  \quad  6 &  1296 & \\
\hline\end{tabular}
\end{table}
\end{document}  